% ============================================================
% 第 8 章  最大熵原理与最小鉴别信息原理
% 第 9 章  组合信息与算法信息
% 教材：石峰、莫忠息《信息论基础》（第三版）P180--P213
% ============================================================

\setcounter{section}{7}
\section{第 8 章\quad 最大熵原理与最小鉴别信息原理}

前七章的 Shannon 信息论都假定概率分布已知；本章处理只掌握部分约束信息
时的推断：Jaynes（1957）\textbf{最大熵原理}与 Kullback（1959）
\textbf{最小鉴别信息原理}。

\subsection{8.1\quad 最大熵原理的提出与含义}

设 $X$ 取值于 $K$ 元集合 $A$，不知 $p(a_k)=P(X=a_k)$，只测得若干函数的期望
$\sum_{k}\p{a_k}f_m(a_k)=C_m$。这组约束一般\textbf{不足以}确定分布（可行解有
无穷多个），故为\textbf{非适定问题}，须补充原则才能选定一个解。

\begin{defbox}{最大熵原理（Jaynes）}
\emph{在只掌握部分信息的情况下对分布作推断时，应取符合约束条件但熵值
最大的分布}——这是唯一"不偏不倚"的选择，任何其他选择都意味着额外添加了
约束或假设。即约束优化问题
\begin{align*}
\max_{P}\ \Hx{X}=-\sum_{k=1}^{K}\p{a_k}\lb\p{a_k}\quad
\text{s.t.}\ \sum_{k}\p{a_k}f_m(a_k)=C_m,\ \
\sum_{k}\p{a_k}=1,\ \p{a_k}\ge0 .
\end{align*}
\end{defbox}

\textbf{合理性（熵集中定理）}：独立重复 $N$ 次试验、第 $k$ 种结果出现
$N_k=Nf_k$ 次，则具该频率的序列数 $w=N!/\prod_kN_k!$；由 Stirling 公式
$n!\sim\sqrt{2\pi n}(n/e)^n$ 得 $\frac1N\lb w\approx-\sum_kf_k\lb f_k=H(F)$，
即 $w(F)\approx2^{NH(F)}$，两组频率的序列数之比 $\approx e^{N(H_1-H_2)}$，
故可行解与最大熵解的序列数之比 $<e^{-N\Delta H}$ 随 $N$ 指数下降：熵值远离
最大值的可行解出现的概率极小，最大熵解是最可靠、最"保险"的解。

\subsection{8.2\quad 拉格朗日乘子法与三个典型结论}

令 $F=H(X)-\beta\bigl(\sum_k\p{a_k}-1\bigr)-\sum_m\lambda_m\sum_k\p{a_k}f_m(a_k)$，
由 $\partial F/\partial\p{a_k}=0$（底取 $e$，不影响结果）得
\begin{align*}
-1-\ln\p{a_k}-\beta-\sum_{m=1}^{M}\lambda_mf_m(a_k)=0
\;\Longrightarrow\;
\p{a_k}=\frac{1}{Z}\exp\Bigl(-\sum_{m=1}^{M}\lambda_mf_m(a_k)\Bigr),
\end{align*}
$Z=\sum_k e^{-\sum_m\lambda_mf_m(a_k)}$，$\beta,\lambda_1,\dots,\lambda_M$ 共
$M+1$ 个常数由 $M+1$ 个约束确定。连续情形在 $\int pf_m\,dx=C_m$、
$\int p\,dx=1$ 下求微分熵最大，同理得
$p(x)=\exp\bigl(-\lambda_0-\sum_m\lambda_mf_m(x)\bigr)$。

\begin{thmbox}{三个典型结论（必背）}
\begin{enumerate}[leftmargin=1.5em,itemsep=1pt,topsep=2pt]
\item 只有归一化约束 $\Rightarrow$ \textbf{均匀分布}：$p_k=1/K$，
      $\Hx{X}=\lb K$；$X\in[a,b]$ 时 $p(x)=1/(b-a)$，$\Hc{X}=\lb(b-a)$。
\item 只给均值 $\Rightarrow$ \textbf{几何/指数分布}：取非负整数、$\E X=m$
      时 $p_k=(1-\theta)\theta^{k}$，$\theta=m/(1+m)$（即
      $p_k\propto\theta^{k}$）；$x\ge0$、$\E X=\mu$ 时
      $p(x)=e^{-x/\mu}/\mu$，$\Hc{X}=\log_2(e\mu)$。
\item 给均值与方差 $\Rightarrow$ \textbf{正态分布}：$\E X=\mu$、
      $\operatorname{Var}X=\sigma^{2}$ 时 $p(x)=N(\mu,\sigma^{2})$，
      $\Hc{X}=\tfrac12\lb(2\pi e\sigma^{2})$。
\end{enumerate}
\end{thmbox}

\textbf{例 8.1.1}：约束 $\E X=a_1$、$\E X^2=a_2$（$X$ 取遍实轴），由
$p(x)=\exp(-\lambda_0-\lambda_1x-\lambda_2x^{2})$ 配方得
$p(x)\sim N(\mu=a_1$，$\sigma^{2}=a_2-a_1^{2})$。
\textbf{例 8.1.2}：掷骰子 1000 次、已知点数均值 $4.5$（$\sum_kkf_k=4.5$），
最大熵解为 $(0.0543,0.0788,0.1142,0.1654,0.2398,0.3476)$，
$H_{\max}=1.61358$ bit（$\lb6=2.585$）；且 $95\%$ 的可行解满足
$1.609\le H\le1.61358$、$99.99\%$ 的可行解满足 $1.602\le H\le1.61358$。

\subsection{8.3\quad 鉴别信息（Kullback--Leibler 距离）}

\begin{defbox}{定义}
设 $X$ 的分布在假设 $H_1,H_2$ 下分别为 $P_1,P_2$，定义鉴别信息（又称相对
熵、交叉熵、方向散度、K--L 距离）
\begin{align*}
D(P_1\Vert P_2)=\sum_{i=1}^{K}P_1(a_i)\log\frac{P_1(a_i)}{P_2(a_i)}
\quad\text{或}\quad
D(P_1\Vert P_2)=\int P_1(x)\log\frac{P_1(x)}{P_2(x)}\,dx .
\end{align*}
\textbf{含义}：为鉴别 $H_1,H_2$，在 $H_1$ 下观察 $X$ 所平均得到的倾向于
$H_1$ 的信息量（认识由先验 $P_2$ 变到后验 $P_1$ 的信息量）；且
$D=-\sum_iP_1\log P_2-\Hx{X}_{P_1}$（交叉熵减熵）。两分布间的\textbf{散度}
$J(P_1,P_2)=D(P_1\Vert P_2)+D(P_2\Vert P_1)$。
\end{defbox}

\begin{thmbox}{鉴别信息的性质}
\begin{enumerate}[leftmargin=1.5em,itemsep=1pt,topsep=2pt]
\item \textbf{非负性}：$D\ge0$，当且仅当 $P_1=P_2$ 时为 0（由
      $\ln x\le x-1$ 证）。
\item \textbf{不对称}：一般 $D(P_1\Vert P_2)\neq D(P_2\Vert P_1)$（方向散度）；
      $J$ 对称、非负，但\textbf{不满足三角不等式}，故不是真正的距离函数，
      只是分布差异的一种度量。
\item \textbf{凸性}：$S=\lambda P+(1-\lambda)Q$ 时
      $D(S\Vert R)\le\lambda D(P\Vert R)+(1-\lambda)D(Q\Vert R)$，
      $D(R\Vert S)$ 亦然。
\item \textbf{独立可加}：$P_i(x,y)=g_i(x)h_i(y)$ 时
      $D(P_1\Vert P_2;XY)=D(g_1\Vert g_2;X)+D(h_1\Vert h_2;Y)$；但\textbf{无
      分步可加性} $D(P_1\Vert P_3)\neq D(P_1\Vert P_2)+D(P_2\Vert P_3)$
      （与 Shannon 熵不同）。
\item 联合与条件鉴别信息 $D(p_1\Vert p_2;XY)=D(g_1\Vert g_2;X)+
      D(g_1\Vert g_2;Y\mid X)$；Johnson（1979）证明满足有限性等条件的泛函
      必为 $B\,D(P_1\Vert P_2)+C\,D(P_2\Vert P_1)$（$B,C\ge0$）。
      计算时须假定所有概率为正，以免除数为 0。
\end{enumerate}
\end{thmbox}

\subsection{8.4\quad 最小鉴别信息原理及其与最大熵原理的关系}

\begin{defbox}{最小鉴别信息原理}
设真实分布密度 $q(x)$ 未知，已知先验密度 $p(x)$ 与约束
$\int qf_m\,dx=C_m$、$\int q\,dx=1$，则应取满足全部约束且使
$D(q\Vert p)=\int q\log(q/p)\,dx$ \textbf{最小}的 $q$ 作为估计——由 $p$ 变到
$q$ 所需信息量最少，$q$ 是约束下最接近 $p$ 的分布。$D(q\Vert p)$ 是 $q$ 的
泛函，属变分问题：按欧拉方程令被积函数对 $q$ 的偏导为零，解得
$q(x)=p(x)\exp\bigl(-\lambda_0-\sum_m\lambda_mf_m(x)\bigr)$，其中 $M+1$ 个
常数由 $M+1$ 个约束确定。
\end{defbox}

\textbf{两原理的关系（高频考点）}：$X$ 为离散分布、先验为等概率
$p(a_k)=1/K$ 时，上式退化为
$q(a_k)\propto\exp\bigl(-\sum_m\lambda_mf_m(a_k)\bigr)$，与最大熵解完全一致：
\emph{最小鉴别信息原理是最大熵原理的推广}（最大熵 = 均匀先验下的最小鉴别
信息）。由 $\ln x\le x-1$ 可证（习题 8.4）：$p_1$ 为最大熵分布、$p_2$ 为任一
满足同一约束的分布时
$\Hx{X}_{p_1}-\Hx{X}_{p_2}=D(p_2\Vert p_1)\ge0$；习题 8.7 给出
$D(g\Vert p)\ge D(g\Vert r)+D(r\Vert p)$（$r$ 为使 $D(r\Vert p)$ 最小的约束内
分布）。

\subsection{8.5\quad 典型应用}

\textbf{(1) 最大熵谱估计}。零均值平稳序列的自相关为
$R(k)=\E(X_tX_{t+k})$，其 Fourier 变换即功率谱
$S(\omega)=\sum_{m}R(m)e^{-im\omega}$（$|\omega|\le\pi$）。周期图法在
$N\to\infty$ 时一般不收敛到真值，Burg（1967）改为按最大熵原则外推自相关。
\textbf{定理 8.1.1（Burg）}：在
$\E(X_tX_{t+k})=\sigma_k$（$k=0,1,\dots$）约束下，熵率最大的序列是 $p$ 阶
自回归高斯过程 $X_t=\sum_{k=1}^{p}\alpha_kX_{t-k}+Z_t$，$Z_t\sim N(0,\sigma^{2})$
且相互独立同分布。用 $X_{t-L}$ 乘两边取均值得
\begin{align*}
R(0)=\sum_{k=1}^{p}\alpha_kR(-k)+\sigma^{2},\qquad
R(L)=\sum_{k=1}^{p}\alpha_kR(L-k),\ L\ge1,
\end{align*}
即 \textbf{Yule--Walker 方程}（Levinson、Durbin 快速算法），由此可外推 $L>p$
的 $R(L)$ 并得谱估计 $S(\omega)=\sigma^{2}/|1+\sum_k\alpha_ke^{-ik\omega}|^{2}$：
在不牺牲分辨率的条件下充分利用了已知知识。\\
\textbf{(2) 图像复原（重建）}：与谱估计并列，是最大熵原理的主要应用领域；
气体分子速度分布（能量约束）、大气密度随高度分布（平均势能约束）也都是
最大熵分布。\\
\textbf{(3) 独立分量分析（ICA，盲分析）}：$Y(t)=AS(t)+Z(t)$，$A$ 未知，仅知
各源独立，以鉴别信息 $D\bigl(p_Y\Vert\prod_ip_{Y_i}\bigr)$ 为优化准则（凸函数
$\Rightarrow$ 全局最优）。

% ============================================================
\section{第 9 章\quad 组合信息与算法信息}

Shannon 理论以"随机模型 + 参数已知"为前提，实际中信源统计特性往往未知
（甚至信号并非随机）。Kolmogorov（1965）指出除统计信息外还有两种定义：
\textbf{组合信息}（计数/枚举）与\textbf{算法信息}（计算机算法），分别导出
\textbf{Fitingof 通用编码}与 \textbf{Lempel--Ziv 通用编码}。

\subsection{9.1\quad 自适应统计编码}

以统计特性为依据的 Huffman 编码称为统计编码，其问题是统计特性不确定
（或信源不平稳）造成的\textbf{统计失配}。设真实分布 $g(a_k)$、码书按
$p(a_k)$ 设计（$-\log_2p$ 为整数），则 $L_{\rm opt}=\Hx{X}$、实际
$\bar L=-\sum_kg\log_2p$、按真实分布的最优码长
$\bar L^{*}=-\sum_kg\log_2g$，三者满足
\begin{align*}
\bar L-\bar L^{*}=\sum_{k}g(a_k)\log_2\frac{g(a_k)}{p(a_k)}=D(g\Vert p)\ \ge0,
\end{align*}
即\textbf{失配的代价恰为鉴别信息}（非负、下凸，仅 $g=p$ 时为 0）。三类编码：
\textbf{静态}编码用固定统计特性；\textbf{半自适应}编码扫描两遍（第一遍统计
造码书、第二遍编码），缺点是须传码书、两次扫描延时大、码书整体更新效率低；
\textbf{自适应}编码只扫一遍，边统计边编码，每步依据"以往消息"的统计特性，
因收发两端历史相同可各自独立更新码书，核心是便于\emph{逐步更新码书}的数据
结构（Huffman、算术编码均可）。自适应编码只略差于最佳静态编码，而静态编码
在失配时远劣于自适应编码。

\subsection{9.2\quad 组合信息}

\textbf{基于组合的信息度量}：设 $X$ 取值于 $N$ 元集合，\emph{不论分布如何}
定义组合信息熵 $H_c(X)=\log_2N$、$I_c=\log_2N$，并平行定义
\begin{align*}
\Hc{X_1\cdots X_K}&=\sum_{i=1}^{K}\log_2N_i,\qquad
\Hc{Y\mid X=a}=\log_2N(Y\mid a),\\
\I{X;Y}_c&=\Hc{Y}-\Hc{Y\mid X},
\end{align*}
$N(Y\mid a)$ 为 $X=a$ 时 $Y$ 取值集合的大小。例如 $X,Y$ 各取 $1,2,3,4$（见
教材表 9.1），$\Hc{Y}=2$，故 $I_c(X{=}1;Y)=0$、$I_c(X{=}2;Y)=2-\log_23$、
$I_c(X{=}3;Y)=1$、$I_c(X{=}4;Y)=2$。组合信息与统计信息性质不同，
\textbf{不可套用} Shannon 熵的公式。

\textbf{组合数与典型序列}：设字母表大小 $K$、序列长 $N$、第 $k$ 个字母出现
$N_k$ 次、频率 $f_k=N_k/N$，则
\begin{align*}
W=\frac{N!}{\prod_{k=1}^{K}N_k!},\qquad
\log_2W\approx N\Bigl(-\sum_{k=1}^{K}f_k\lb f_k\Bigr)=N\,H(f).
\end{align*}
二元情形 $W=C(N,k)$，取 $k=\alpha N$ 得
$\log_2C(N,\alpha N)\approx NH_2(\alpha)$，且
$2^{NH_2(\alpha)}/(N+1)\le C(N,\alpha N)\le2^{NH_2(\alpha)}$，其中
$H_2(\alpha)=-\alpha\lb\alpha-(1-\alpha)\lb(1-\alpha)$。\textbf{"用典型序列
所占比例定义信息量"}：长 $N$ 的二元序列共 $2^{N}$ 个，"恰有 $k=\alpha N$ 个
1"这一类（典型序列组）约有 $2^{NH_2(\alpha)}$ 个，占全体的比例
$\rho=C(N,k)/2^{N}\approx2^{-N[1-H_2(\alpha)]}$，故
$-\log_2\rho\approx N[1-H_2(\alpha)]$；
描述"某序列属于哪一组"只需 $\log_2C(N,k)\approx nH(k/n)$ bit，再加约
$\log_2(N+1)$ bit 指明 $k$（定理 9.3.4 的思路）。信息量由\emph{可能的组合
个数}决定、与分布无关——这就是组合信息的实质。

\subsection{9.3\quad Fitingof 通用编码}

思想：不用统计特性，只用"计数 + 序号"。若某组序列的频率使
$\sum_kf_k\log_2f_k$ 在常数 $h$ 附近，则序列数 $W\approx2^{Nh}$，用 $Nh$ 位
二元码即可编码。步骤：(1) 序列分成长为 $N$ 的组（$N$ 充分大）；(2) 统计频率
分布 $F$，不同频率分布总数为 $W(K,N)$，对 $F$ 编码需
$\lceil\log_2W(K,N)\rceil$ bit；(3) 同一 $F$ 下仍有 $W(F)$ 个序列，区别它们需
$\lceil\log_2W(F)\rceil$ bit。故
\begin{align*}
L=\Bigl\lceil\log_2W(F)\Bigr\rceil+\Bigl\lceil\log_2W(K,N)\Bigr\rceil,\quad
W(F)=\frac{N!}{\prod_kN_k!},\ \
W(K,N)=C_{N+K-1}^{N}.
\end{align*}
\textbf{定理 9.2.1（通用性）}：$\forall\varepsilon>0$、$\exists N_0$，当
$N>N_0$ 时 $\frac1N\E\bigl[L(X_1\cdots X_N)\bigr]<H(P)+\varepsilon$，即对不同
统计特性的信源都实现理想压缩（证明用 Stirling 公式放缩
$\log_2C_{N+K-1}^{N}$，并用 Jensen 不等式得
$\E\log_2\bigl(N!/\prod N_k!\bigr)\le NH(P)$）。该差值的下降速率\emph{缓慢}
（统计编码更快），故有限码长时统计编码性能更好，该法工程上尚少用。

\subsection{9.4\quad 算法信息}

\textbf{Kolmogorov 算法熵}：$S_1=0101\cdots01$（16 个"01"）很简单；$S_2$ 是某
常数的二进展开，看着随机但仍简单；$S_3$ 只能用"1 的个数 + 序号"描述，复杂性
与长度成正比。规律越强程序越短 $\Rightarrow$ 程序长度可度量信息量。

\begin{defbox}{定义 9.3.1}
设 $p$ 是给定 $y$ 时使计算机 $C$ 输出 $x$ 的二元程序、$l(p)$ 为其长度，称
$K(x\mid y)=\min_{p:\,C(p,y)=x}l(p)$ 为 $x$ 在条件 $y$ 下的\textbf{条件算法
熵}；$K(x)=K(x\mid\varnothing)$；$I(x;y)=K(x)-K(x\mid y)$。任一通用 Turing 机
上的程序可在另一台上执行（加前缀仿真程序），两台机器算得的熵只差一个与
$x,y$ 无关的常数，故"计算机"理解为某台通用 Turing 机即可。
\end{defbox}

性质：(1) $K(x\mid l(x))\le l(x)+C$；(2)
$K(x)\le K(x\mid l(x))+2\log_2l(x)+C$（长度本身也要传递）；
(3) $|\{x:K(x)<k\}|<2^{k}$——长度小于 $k$ 的程序只有
$1+2+\cdots+2^{k-1}<2^{k}$ 个，故\textbf{低复杂性串极少}；(4)
$K(x_1\cdots x_N\mid N)\le\log_2N+NH(k/N)+C$（$k=\sum_ix_i$，给定 $k$ 与序号
即可复现）；(5) $NH(X)\le\E[K(x\mid N)]\le NH(X)+\log_2N+C$（左不等式由停机
程序集合构成前缀码、满足 Kraft 不等式导出）；(6) 绝大多数序列满足
$\lim_{N\to\infty}K(x_1\cdots x_N)/N=1$，即\textbf{不可压缩序列}（真正的随机
序列，算法熵故又称\textbf{算法随机度}）。注意 Shannon 的 $I(X;Y)=I(Y;X)$、
$H(X,Y)=H(X)+H(X\mid Y)$ 在算法信息论中一般不成等式，只在渐近意义下近似成立。

\textbf{例 9.3.1}：$e$ 的平方根的前 $31239872948332$ 位比特构成极长二元串，
用一句话即可精确描述，共 70 个字符；按 8 bit ASCII 计仅
$8\times70=560$ bit，故其 Kolmogorov 复杂性不超过 560 bit。\\
\textbf{不可计算性}：1970 年 Chaitin 证明不存在一般办法算出任一串的算法熵。
直观论证：若可计算，则"打印出 $K(x)\ge10^{10}$ 的第一个二元序列"本身就是很短
的描述，导出矛盾（与 Perry 悖论、Gödel 不完备定理同源）。故该理论的价值在于
研究随机性与判定问题，而非求最短程序。

\textbf{Lempel--Ziv 通用编码}：理想通用编码不可实现，于是限制计算模型——Ziv
与 Lempel（1978）基于有限状态自动机提出有限状态压缩编码与有限状态熵。把序列
迭代分解为\emph{互不相同}的最短新短语 $w_1,w_2,\dots$（第 $i$ 步从上一短语后
第一个字母起，找此前未出现过的最短短语）；$w_i$ 去末字母后的前缀必已出现过，
若它第 $j<i$ 步出现则用 $(j,c_i)$ 表示 $w_i$。$j$ 需 $\lceil\log_2i\rceil$ bit、
末字母需 $\lceil\log_2M\rceil$ bit（$M$ 为源字母表大小），短语互不相同故分解
唯一、无需分隔符：
\begin{align*}
L=\sum_{i=1}^{d}\bigl(\lceil\log_2i\rceil+\lceil\log_2M\rceil\bigr),
\qquad
\limsup_{N\to\infty}\rho(x^{N})\le H(X)+\varepsilon,
\end{align*}
$\rho=L/(N\log_2M)$ 为压缩比。码率渐近于熵率 $H(X)$ 且不用统计特性
$\Rightarrow$ 通用；LZ 已成为文件压缩标准（UNIX 标准压缩程序、PC 的 arc），
常把 ASCII 文本压到约一半。

\textbf{Kieffer--Yang 通用编码}（2000）：\textbf{基于语法}的通用编码——用重复
出现的子串建立语法短语 $S_1,S_2,\dots$。例如对 $X=00110011100111001001\cdots$
随读入把重复子串定义为新短语（如 $S_3=00$、$S_2=S_11$ 等），最终把整段写成
少数短语的级联 $X=S_1S_2S_2S_1S_3\cdots$；传输时只需短语定义与短语号码串即
可无失真复制（分步例见教材 P208--P209）。它把 LZ 的字符串匹配与算术码结构
结合，优于 LZ 与算术码；计算复杂度是序列长度的线性函数；平稳遍历时码率趋于
$H(X)$。

\subsection{9.5\quad 近似熵与样本熵}

近似熵（ApEn，Pincus 1991）从衡量时间序列复杂性的角度提出，已成功用于心率、
血压等生理序列。记第 $i$ 个 $m$ 长片段 $x^{m}(i)=(u_i,\dots,u_{i+m-1})$。

\begin{defbox}{ApEn 的计算与定义}
(1) 片段距离 $d\bigl(x^{m}(i),x^{m}(j)\bigr)
=\max_{0\le k\le m-1}|u_{i+k}-u_{j+k}|$；(2) 给定阈值 $r$，与 $x^m(i)$ 距离
$<r$ 的片段数（模板匹配数）记为 $N^{m}(i)$，比重
$B^{m}(i)=N^{m}(i)/(N-m)$；(3) 取对数平均
$\varphi^{m}(r)=\frac{1}{N-m+1}\sum_{i=1}^{N-m+1}\ln B^{m}(i)$；(4) 对 $m+1$
重复得 $\varphi^{m+1}(r)$，则
\begin{align*}
\mathrm{ApEn}(m,r,N)=\varphi^{m}(r)-\varphi^{m+1}(r).
\end{align*}
\end{defbox}

\textbf{含义与用途}：近似熵相当于维数由 $m$ 增到 $m+1$ 时\emph{新模式出现的
对数条件概率的均值}，反映序列结构上的复杂性，对随机过程与确定性过程都适用：
近似熵越大 $\Rightarrow$ 产生新模式概率越大 $\Rightarrow$ 序列越复杂（周期
信号很小、随机信号较大）。它与幅值、相位无关，只与频率成分的复杂性有关，
便于分析实测数据；缺点是结果不恒定，严重依赖序列长度与弱的自相似性。

\textbf{样本熵}（SampEn，Richman 等 2000）：记号同前，但只在\textbf{其余}
$N-m$ 个片段中统计匹配数 $B^{m}(i)=N^{m}(i)/(N-m)$，对 $i$ 取平均得
$B^{m}(r)$，对 $m+1$ 同法得 $B^{m+1}(r)$，则
\begin{align*}
\mathrm{SampEn}(m,r,N)=-\ln\frac{B^{m+1}(r)}{B^{m}(r)} .
\end{align*}
\textbf{异同（必考）}：物理意义相近——样本熵越低，序列自我相似性越高，越高则
越复杂。关键差别：近似熵把数据与\emph{自身}比较（自身匹配计入）以避免
$\ln0$，但引入偏差；样本熵\textbf{不计自身匹配}，因为熵是新信息产生率的测度，
比较数据与其自身毫无意义，故精度更高、一致性更好，更适于特征提取。二者都
依赖 $m,r$：$m$ 一般取 $2\sim10$，维数由数据结构定；$r$ 亦依数据而定，尚无公认
最佳标准（教材 P212--P213 此处 OCR 有缺损，以教材原文为准）。

\subsection{9.6\quad 本章作业与讲评：习题 9.2}

\begin{example}{题目（教材 P213，习题 9.2）}
试对字母序列 \texttt{ababbbabaabababaababbababa} 进行 Lewpel--Ziv 编码。
（教材拼作 Lewpel--Ziv，即 Lempel--Ziv；序列共 26 个字母，$a$、$b$ 各 13 个。）
\end{example}

\textbf{思路}：(1) 从左到右把序列切成互不相同的短语——在已建立的短语字典中
找当前最长的可匹配短语 $w_j$，再附上紧随其后的一位字母，即得新短语 $w_i$；
开头字母无匹配前缀，直接编码为 $(0,c)$。这样每个短语的前缀必是更早出现的
短语，$j$ 才有定义、分解才唯一。(2) 码字：$w_i=w_jc_i$ 记为
$\bigl(j$ 的二进制$,\,c_i\bigr)$，$j$ 用 $\lceil\log_2i\rceil$ 位（$j=0$ 时
省去），字母按 $a\to0$、$b\to1$ 用 1 位；(3) 校验：短语长度之和等于 26，
短语两两不同故无需分隔符，可无失真重构。

\begin{keypoint}{解答（LZ 字典法分解与码字）}
分解 $a\mid b\mid ab\mid bb\mid aba\mid abab\mid abaa\mid ba\mid bba\mid
bab\mid a$，长度 $1{+}1{+}2{+}2{+}3{+}4{+}4{+}2{+}3{+}3{+}1=26$。码字依次
$a\!\to\!(0,0)$、$b\!\to\!(0,1)$、$ab\!\to\!(01,1)$、$bb\!\to\!(10,1)$、
$aba\!\to\!(011,0)$、$abab\!\to\!(101,1)$、$abaa\!\to\!(101,0)$、
$ba\!\to\!(010,0)$、$bba\!\to\!(0100,0)$、$bab\!\to\!(1000,1)$、
$a\!\to\!(0001,\varnothing)$（末短语恰为已有短语 $w_1=a$，故不附新字母）。
前 6 个码字与教材 P208 示例
（$a.b.ab.bb.aba.abab\cdots\to(0,0)(0,1)(01,1)(10,1)(011,0)(101,1)$）完全吻合，
可作自检。总码长
$L=\sum_{i=1}^{11}\bigl(\lceil\log_2i\rceil+1\bigr)=1+2+3+3+4+4+4+4+5+5+4
=39$ bit，而 26 个二元字母原只需 26 bit——短序列中短语号码冗余很大，LZ 的
优势体现在长序列的渐近意义。
\end{keypoint}

\begin{warning}{讲评：常见错误}
(1) 短语须两两不同、按"在已出现部分中搜索"切分，否则 $j$ 无意义、接收端
无法复原；(2) $j$ 的位数随 $i$ 增长取 $\lceil\log_2i\rceil$，不是定长；
(3) 须声明字母与比特的约定（$a\to0$、$b\to1$），否则答案不唯一；(4) 本题是
LZ 编码，勿与 Huffman、算术编码混淆——若课件按 LZ77 三元组（偏移量、匹配
长度、下一字符）讲，也可写成 $a\mid b\mid abb\mid baba\mid ababa\mid
baababb\mid ababa$，对应
$(0,0,a)(0,0,b)(2,2,b)(4,3,a)(9,4,a)(7,6,b)(10,5,\varnothing)$；两种写法都对，
但须与课上要求的形式一致。
\end{warning}
